\section{Flat Nil paths on the fixed compact quotient}
\label{app:nilquotient}

The paths $T_2,T_3$ lie in the positive determinant-one horizontal metrics
below, with the fixed coordinate periods and gluing of Section 2. Their shape
variable $(M^{\mathsf T}M)_{22}$ remains one, so (H-03) is constant along them.

\begin{equation}\tag{N-01}\label{eq:N-01}
\begin{aligned}
g_Q&=r^2\left[(\mathrm{d}x,\mathrm{d}w)Q(\mathrm{d}x,\mathrm{d}w)^{\mathsf T}
+(\mathrm{d}z-x\mathrm{d}w)^2\right]+L^2\mathrm{d}y^2,\\
Q&>0,\qquad\det Q=1,\qquad
G_{H^1}=\frac{\mathcal V}{4\pi^2r^2}Q^{-1}.
\end{aligned}
\end{equation}

The last expression is the $L^2$ Gram matrix of the integral one-forms
$\mathrm{d}x/(2\pi),\mathrm{d}w/(2\pi)$, with $\mathcal V=(2\pi)^4Lr^3$.
These forms are closed and coclosed: the horizontal block of the inverse
coordinate metric is $r^{-2}Q^{-1}$ and the volume density is constant, so
their divergences vanish. The abelianization of the Nil lattice has rank two;
these forms, together with $\mathrm{d}y$ on the product, span harmonic degree one.
The LC Ricci eigenvalues are $-1/(2r^2)$ horizontally, $+1/(2r^2)$ centrally,
and zero on the circle. Any product-metric isometry therefore preserves the
horizontal harmonic subspace and its integral lattice. At fixed $r,L$, its
normalized shortest squared norm is an isometry invariant:

\begin{equation}\tag{N-02}\label{eq:N-02}
\begin{aligned}
\sigma(Q)&=\min_{m\in\mathbb Z^2\setminus\{0\}}m^{\mathsf T}Q^{-1}m,
\qquad a=\sqrt2q,\\
Q_{T_2}&=\mathrm{diag}(e^a,e^{-a}),
&\sigma(Q_{T_2})&=e^{-|a|},\\
Q_{T_3}&=\begin{pmatrix}\cosh a&\sinh a\\\sinh a&\cosh a\end{pmatrix},
&\sigma(Q_{T_3})&=\min(\cosh a,2e^{-|a|}).
\end{aligned}
\end{equation}

For $T_2$, the smaller diagonal coefficient bounds every nonzero integer
vector and is attained by a coordinate vector. For $T_3$, integer vectors of
Euclidean squared length one give $\cosh a$; all others have squared length
at least two and hence norm at least $2e^{-|a|}$, attained by $(1,\pm1)$
with the appropriate sign. This proves (N-02) on the infinite lattice.

Distinct $|q|$ are distinguished on the entire $T_2$ path and on the interval
$0\le\sqrt2|q|<\log(3)/2$ of $T_3$. In that interval, nonzero $T_2$ points
have $\sigma<1$ and nonzero $T_3$ points have $\sigma>1$, so the paths are
also distinguished from each other there. Opposite signs are nevertheless
identified by explicit lattice-preserving maps. In the normalized coordinates
$(X,Y,Z)=(x/(2\pi),w/(2\pi),z/(4\pi^2))$, these are

\begin{equation}\tag{N-03}\label{eq:N-03}
(X,Y,Z)\mapsto(Y,-X,Z-XY)\quad(T_2),\qquad
(X,Y,Z)\mapsto(-X,Y,-Z)\quad(T_3).
\end{equation}

They preserve the integer Heisenberg group law and the square of
$\mathrm{d}Z-X\mathrm{d}Y$, while taking $Q(q)$ to $Q(-q)$.
Horizontal determinant-one transformations give local Heisenberg
automorphisms, but need not preserve the lattice. This local/global distinction
is consistent with the geometry of fixed Heisenberg quotients\cite{Boldt2021}.
Here the integral harmonic norm serves as an obstruction to isometry.
The flat paths contain continuously distinct compact metrics, along which
the transverse derivatives in (H-05) remain nonzero.
